\section{The exact mixed tangent image needed by the extraction proof}
\label{sec:cubic-tangent-image}

This section gives a conditional replacement for one precise hypothesis
in the source extraction argument.  It does not construct a cubic
\texttt{Dat}, identify its integral Selmer module, or assert that its
mixed tangent image vanishes.  Paths in this section are relative to
\texttt{src/Development/AutomorphicArithmetic/AutomorphicVarieties/}
at the pinned commit of AC-K29, except where a full relative path is
explicitly written.

\paragraph{The literal source domains.}
Write $O=c.O$ and $\kappa=\overline{\mathbb Q}_\ell$ with the
source coefficient embedding.  In
\path{UnitaryThreeGroups/TanCocIsBK.lean},
\texttt{tanCoc\_isBK}, the quantifier is over every
\[
 F:S.R\longrightarrow_O\kappa[\epsilon]/(\epsilon^2),
 \qquad \operatorname{IsLocPt}(\operatorname{dualFst}F).
                                                               \tag{GT.1}
\]
Here $S:D.\operatorname{Ring}(U,g_1)$, $U:D.\operatorname{ULoc}$,
$D.\operatorname{Marked}(g_1)$, and $D:\operatorname{Dat}
(\Omega,K,c,e,\psi)$ have exactly the source meanings.
In particular $\ell$ is prime, $\Omega$ is an algebraic closure
of $\mathbb Q$, $K$ is a number field embedded in $\Omega$,
$N\ne0$, and $\psi$ is the source Dirichlet character modulo $N$.
The theorem also assumes its entire \texttt{CohInp}.
It proves that the displayed cocycle is continuous and satisfies
all of the source's \texttt{IsBK} local conditions; it does not
prove that its global class is zero.

The relevant restriction on the specialized point is in
\path{UnitaryThreeGroups/Rigid0.lean}, record \texttt{PtHyp}.
Its last field is
\[
 \forall f:S.R\longrightarrow_O\kappa,\quad
 \operatorname{TwoPt}(\eta,r,f)\Longrightarrow
 \operatorname{Rigid}(r,f),\qquad r=S.r.\operatorname{hom}.
                                                               \tag{GT.2}
\]
In \path{UnitaryThreeGroups/GRepresents.lean},
\texttt{tanCoc}, \texttt{Rigid}, and \texttt{RigidWt} are defined explicitly.
The predicate \texttt{Rigid} quantifies over every $F$ in (GT.1)
with $\operatorname{dualFst}F=f$, not just lifts chosen by one
basis of a tangent space.  Its conclusion is
$\operatorname{tanCoc}(F,r)=\operatorname{cob}(r_f,X)$ for
some matrix $X$, where $r_f=\operatorname{push}(f,r)$.
Here \texttt{TwoPt}, defined in \texttt{ConjStable.lean}, means
\texttt{IsLocPt} together with neither \texttt{Upper} nor
\texttt{Lower} for the specialized representation.
The other two fields of \texttt{PtHyp} are retained: at every local
point which is both upper and lower, its specialized representation
equals the base representation (not an asserted equality of algebra
maps), and
\texttt{RigidWt} holds at that base point for every integer
weight $j\le0$.

The statement of
\path{UnitaryThreeGroups/ModelUpper.lean}, \texttt{model\_upper},
is more algebraic.  Its quantifiers range over every prime
$\ell\ne2$, coefficient $c$, topological groups $\Gamma,H$,
characters $\eta:\Gamma\to\{\pm1\}$ and
$\Psi:\Gamma\to O^\times$, set $\mathcal N\subset\Gamma$,
continuous homomorphism $\iota:H\to\Gamma$, the specified universal
local ring $U$ for the restrictions of $(\eta,\Psi)$, an ideal
$I\subset U.R$, elements $c_c,g_1\in\Gamma$, a source
$S:\operatorname{SRing}(\mathcal N,\iota,U,I,c_c,g_1)$, and every
$M:\operatorname{Model}(S,\kappa)$.
It retains \texttt{Marks}, \texttt{LemmaM}, finiteness over $O$
of the weight-zero subalgebra, \texttt{PtHyp} at $M.x_0$, and
nonvanishing of the selected trace difference at every two-point
point.  None of those conditions is dropped below.
The proof uses the last field of \texttt{PtHyp} in
\path{UnitaryThreeGroups/GradedModelDerivationConditionOfPtHyp.lean},
\texttt{gradedModel\_derivationCondition\_of\_ptHyp}, to control
all point derivations, before producing the finite trace family.

\begin{proposition}[AC-K29-tangent: exact replacement of the mixed rigidity step]
\label{AC-K29-tangent}
Fix any of the source data just described, or a proposed modified
arithmetic model for which the same ring, lift, local-condition and
extraction interfaces have separately been established.
The following assertions are pointwise in all of those data; to
apply them throughout the family in AC-K29 they must hold for every
$x,e,\ell,o$ and every coefficient extension, local choice $U$,
marked element $g_1$, ring $S$, model $M$, two-point point $f$, and
lift $F$ used in that extraction.

Suppose each type-II point has endpoint block $B_f$, middle
character $m_f$, and a normalized semidirect unitary action on
$\operatorname{Mat}_3(\kappa)$ preserving the decomposition
\[
 \operatorname{ad}(r_f)=\mathcal A_f\oplus\mathcal C_f
                                      \oplus\mathcal R_f.       \tag{GT.3}
\]
Here $\mathcal A_f$ consists of the four entries in the endpoint
corner $\{0,2\}\times\{0,2\}$, $\mathcal C_f$ of the $(1,1)$
entry, and $\mathcal R_f$ of the four endpoint/middle entries.
Denote the coordinate projections by $P_A,P_C,P_R$.  Thus
\[
 P_RX=\begin{pmatrix}
 0&X_{01}&0\\ X_{10}&0&X_{12}\\0&X_{21}&0
 \end{pmatrix}.                                               \tag{GT.4}
\]
Assume, for every such $f$ and every $F$ with
$\operatorname{dualFst}F=f$, that the cocycle
$\gamma_F=\operatorname{tanCoc}(F,r)$ is continuous and satisfies
\texttt{IsBK}, and that $P_A\gamma_F$ and $P_C\gamma_F$ are
coboundaries in their respective representations.
Define the actual mixed tangent image by
\[
 \mathcal T_f=
 \left\{[P_R\gamma_F]:
 F\in\operatorname{Hom}_{O\text{-alg}}
       (S.R,\kappa[\epsilon]/(\epsilon^2)),\
       \operatorname{dualFst}F=f\right\}
       \subseteq H^1_f(\mathbb Q,\mathcal R_f),                 \tag{GT.5}
\]
when $\Gamma=G_{\mathbb Q}$; for a general group use the
continuous $H^1$ with the stipulated local conditions instead.
Then
\[
 \mathcal T_f=\{0\}
 \quad\Longleftrightarrow\quad
 \forall F\text{ over }f,\quad
       \gamma_F\text{ is a coboundary}.                         \tag{GT.6}
\]
The assertion on the right is exactly the type-II part of (GT.2).
It is also equivalent to the mixed cocycle of every actual lift
being a coboundary, under the retained diagonal-block assumptions.

Consequently, if the first two fields of \texttt{PtHyp}, the
type-I/type-II dichotomy, the entire type-I rigidity step,
and all other premises and integral
comparisons of \texttt{model\_upper} and \texttt{vlen\_le} are
retained, imposing $\mathcal T_f=0$ at every type-II point supplies
the remaining rigidity field.  Together with the actual automorphic
admissibility, trace-class injection, and local-condition
identification, this is sufficient for the extraction premise of
AC-K29.  It is not a proof of those additional premises.

For the cubic endpoint descent of AC-G45, before or after the
fixed scalar normalization of AC-G46, the mixed module is
\[
 \mathcal R_f\simeq V_f\otimes
                  \operatorname{Ind}_{G_K}^{G_{\mathbb Q}}\chi^j
       \quad (j=1,2).                                        \tag{GT.7}
\]
The full assertion $H^1_f(\mathbb Q,\mathcal R_f)=0$ is sufficient
for (GT.5) to vanish.  The replacement (GT.5) asks only about the
image of the actual lifts.  No surjectivity of this image onto
$H^1_f$ is assumed or proved.

Finally, a fixed scalar character twist has the exact identities
\[
 \operatorname{ad}(r_f^\nu)=\operatorname{ad}(r_f),\qquad
 \operatorname{tanCoc}(F,r^\nu)=\operatorname{tanCoc}(F,r).
                                                               \tag{GT.8}
\]
They hold also for the anticyclotomic semidirect extension of
$\nu$ in AC-G46.  Thus a transport by that fixed twist leaves
the image (GT.5) unchanged whenever the actual rings and lifts
are transported; it does not by itself prove its vanishing.
\end{proposition}
\begin{proof}
\proofstep{1. The actual tangent cocycle.}
Write $F(a)=f(a)+\epsilon d(a)$, so that $d$ is an $O$-derivation
with values in the $S.R$-module $\kappa_f$.  If $a(g)$ is the
matrix component of $r(g)$, the literal definition is
\[
       \gamma_F(g)=d(a(g))\,f(a(g))^{-1}.                       \tag{GT.9}
\]
The first-order product rule, including the derivative of
$J a^{-t}J$ on the nonidentity component, gives
$\gamma_F(gh)=\gamma_F(g)+\operatorname{ad}(r_f)(g)\gamma_F(h)$.
This is exactly \texttt{tan\_coc} and its complete matrix proof
in \texttt{TanCocIsBK.lean}.  The continuity and local conditions
are the retained hypotheses, supplied by \texttt{tanCoc\_isBK}
for the literal source data.  Linear combinations of $O$-derivations
again give lifts of $f$, and the zero derivation gives the constant
lift.  Thus (GT.5) is a linear subspace containing zero.

\proofstep{2. The source block decomposition.}
For a block matrix with endpoint corner $B_f$ and middle scalar
$m_f$, conjugation preserves each of the three coordinate spaces
in (GT.3).  The other component of the semidirect action is
negative transpose followed by the anti-diagonal reversal;
it preserves the endpoint corner and the middle entry and
exchanges the two mixed rectangular spaces.  Hence all three
coordinate projections commute with the full adjoint action.
They are continuous, sum to the identity, and have pairwise
zero products.

For comparison with the literal old source, let
\[
 p_\pm(X)=(X_{01}\pm X_{12},\ X_{21}\mp X_{10}),\qquad
 s_\pm(v)=\begin{pmatrix}
 0&v_0&0\\ \mp v_1&0&\pm v_0\\0&v_1&0
 \end{pmatrix}.
                                                               \tag{GT.10}
\]
A direct entrywise calculation gives
$P_RX=\tfrac12s_+p_+(X)+\tfrac12s_-p_-(X)$.
These are exactly \texttt{pB 1}, \texttt{pB (-1)},
\texttt{sB 1}, \texttt{sB (-1)} and \texttt{decomp} in
\path{TheoremsMuLinked/GaloisRepresentationTypes/ReducibleThreeDim.lean}.
In the old determinant-one endpoint situation their equivariance
identifies the two summands with $V_f$ and $V_f\eta_K$.
In the cubic situation the whole mixed space remains invariant,
whereas these two old summands need not be invariant separately.
The coset-column identification proved in AC-G45 gives (GT.7).
The combined projection (GT.4), rather than an unproved old
rank-two splitting, is therefore the relevant one.

Continuous equivariant projections preserve cocycles,
coboundaries and the stated local conditions.  For the latter,
away from $\ell$ apply the projection to an inertia coboundary;
at $\ell$ apply its tensor with $B_{\mathrm{cris}}$ to the
trivializing vector.  This is also the complete argument of
\texttt{bk\_map} in the same source file.  Thus (GT.5) has the
asserted target.

\proofstep{3. Equivalence and the exact source use.}
If $[P_R\gamma_F]=0$, choose $Y_R\in\mathcal R_f$ with
$P_R\gamma_F(g)=\operatorname{ad}(r_f)(g)Y_R-Y_R$.
Choose $Y_A,Y_C$ similarly using the retained assumptions.
Their sum gives
\[
 \gamma_F(g)=\operatorname{ad}(r_f)(g)(Y_A+Y_C+Y_R)
                                      -(Y_A+Y_C+Y_R).
\]
The source convention is
$\operatorname{cob}(r_f,X)(g)=X-\operatorname{ad}(r_f)(g)X$;
take $X=-(Y_A+Y_C+Y_R)$.
Conversely, projecting any such coboundary proves that its mixed
class is zero.  This proves both directions of (GT.6), with
exactly the stated quantifier over all lifts of each point.

The entire proof of
\path{ModularityArgument/ThresholdArguments/DenseThresholdCases/FPtHypU.lean}
first obtains \texttt{tanCoc\_isBK}.  At a two-point point it
uses type-I rigidity in one case and full adjoint vanishing in
the type-II case, and then converts membership in $B^1$ to the
source \texttt{cob} convention.
The complete proof of \texttt{TypeIIOfBlochKatoLocalPoint.lean}
in the same directory handles the corner with
\texttt{RigiditySecondCase}, the middle with the nonpositive-weight
input, and the old two mixed summands with \texttt{kUse'}.
Replacing only that last use by (GT.5) and the preceding
recombination gives precisely the same last field of
\texttt{PtHyp}.  All first-field, base-weight, type-I,
integrality, finite-algebra, model and admissibility hypotheses
remain present.  The complete \texttt{model\_upper} proof then
applies to literal source models; a proposed modified model must
also establish its corresponding interface before this inference
can be used for the cubic coefficient.

\proofstep{4. Fixed scalar twisting.}
For a usual representation, write $a^\nu(g)=\nu(g)a(g)$.
Scalar matrices cancel in conjugation, so the adjoint action is
unchanged.  For the semidirect unitary group use a scalar function
$q(g)$ satisfying $q(gh)=q(g)q(h)^{\eta(g)}$, as furnished by
an anticyclotomic character on $G_K$ and value one at complex
conjugation.  Then $a^\nu(g)=q(g)a(g)$ defines the corresponding
semidirect twist; scalars still cancel in
$a^\nu(g)\theta_{\eta(g)}(X)a^\nu(g)^{-1}$.
The fixed twist is valued in the coefficient algebra, so every
$O$-derivation has $d(q(g))=0$.  Equation (GT.9) gives
\[
 d(q(g)a(g))\,f(q(g)a(g))^{-1}
   =q(g)d(a(g))\,f(a(g))^{-1}q(g)^{-1}
   =\gamma_F(g).
\]
This proves (GT.8), including the cocycle itself, not just its
cohomology class.  Coordinate projections and local-condition
cochain complexes are therefore identical.  Transport of a
particular deformation ring or its prescribed local conditions
is a separate hypothesis; none has been inferred from this
identity.
\end{proof}

\paragraph{Admissibility and scope.}
The fixed-$K$ counting class has no explicit determinant-one
condition.  Its \texttt{Adm3Alt} definition is in
\path{ModularityArgument/ThresholdArguments/DenseThresholdCases/AfStmt.lean},
and \texttt{Adm3U\_s} adds \texttt{Unl} in \texttt{NumberU3US.lean}.
The imported definition
\path{src/Development/AutomorphicArithmetic/Foundations/GaloisRepresentations/AdmissibleResidues/IsAutomorphicOverImagQuad.lean},
\texttt{AutCsd0}, asks for an attached regular algebraic cuspidal
representation satisfying \texttt{DualIsConj}; it does not impose
a trivial central character.  The remaining \texttt{WitA},
\texttt{Level} and \texttt{Unl} conditions still require proof.
At the inert primes of AC-K29, the blockwise map
$(\nu\chi^j V_f)|_{G_K}\oplus\nu\mapsto V_f|_{G_K}\oplus1$
preserves the tested trace, and is injective on isomorphism classes
when the endpoint block is irreducible and the twisting characters
are fixed.  Indeed the unique irreducible two-dimensional summand
is recovered from either side and the fixed twist is invertible.
Membership of the image in the source \texttt{Adm2A} still asks
for its actual newform attachment, \texttt{WitAQ}, determinant
weight and level.  Neither this counting observation nor the
scalar determinant correction supplies the vanishing in (GT.5).
